% Source: MQ17a_v1.html

\documentclass[11pt]{article}
\usepackage[margin=1in]{geometry}
\usepackage{amsmath}
\usepackage{amssymb}
\usepackage{siunitx}
\usepackage{enumitem}

\title{MQ17a: Temperature and Thermal Expansion}
\author{}
\date{}

\begin{document}
\maketitle

\section*{MQ17a: Temperature and Thermal Expansion}

\begin{enumerate}[leftmargin=*,label=\textbf{Question \arabic*.}]

\item \hfill \textit{1 pts}

Convert 61°C into Kelvin.

\item \hfill \textit{1 pts}

A certain metal has a coefficient of linear expansion of 4 × 10$^{-5}$ (C°)$^{-1}$.  Suppose a pipe made of this metal is 2.4 m long. Calculate the temperature change in Celsius if the length changes by -4.7 mm. (Indicate a decrease by a negative change in temperature.)

\item \hfill \textit{1 pts}

A substance is cast into a bar 11.5 m long. When the bar's temperature is raised from 5°C to 62°C, its length increases by 0.018 m. Calculate the coefficient of linear expansion for this material and supply the missing coefficient in the blank below.



𝛼 = \rule{2.5cm}{0.4pt} × 10$^{-5}$ (C°)$^{-1}$

\item \hfill \textit{1 pts}

A block of metal with dimensions 2.8 m × 1.2 m × 1.6 m is warmed by 5.9 C°. Calculate its volume increase in cm$^{3}$ if its coefficient of linear expansion is 1.9 × 10$^{-5}$ (C°)$^{-1}$.



Δ\emph{V} = \rule{2.5cm}{0.4pt} cm$^{3}$

\item \hfill \textit{1 pts}

Suppose 4,498 J of energy are transferred to a block of mass 1.5 kg and specific heat 637 J/(kg⋅C°). Calculate the temperature change in Celsius degrees.

\item \hfill \textit{1 pts}

Suppose a cube of iron of mass 6.5 kg and temperature 89°C falls into 0.3 kg of water at 24°C. Neglecting the container and any energy loss, calculate the final equilibrium temperature given \emph{c}$_{Fe}$ = 444 J/(kg⋅C°) and \emph{c}$_{water}$ = 4190 J/(kg⋅C°).



\emph{T} = \rule{2.5cm}{0.4pt} °C

\end{enumerate}

\end{document}
